  Pro opce evropského typu je vyhodnocení v celku jednoduché: Stačí vzít finální cenu aktiva, dle typu opce spočíst výdělek, získat průměr a ten následně slevit o růst ceny měny za daný čas. Nicméně pro opce amerického typu musíme počítat s možností dřívějšího využití opce.

  (Dle prezentace) Nacházení bodu předčasného výkonu opce je antithetické pro Monte Carlo metody, poněvadž pro nalezení bodu předčasného výkonu musíme cestovat od konce k začátku(známe cenu na konci, musíme najít nejlepší hodnotu dané cesty), přičemž Monte Carlo jde od začátku ke konci (generování ceny aktiva pomocí Geometrického Brownovského pohybu).
  Formulace problému nalezení předčasného výkonu spadá pod dynamické programování. Dynamické programování je kategorie problémů (TODO: Dopsat definici).

  V rámci zjednodušení, zaveďme si dvě funkce: $h_i(x)$ a $V_i(x)$, kde $h_i(x)$ označuje hodnotu předčasného výkonu opce $x$ v čase $t_i$, diskontovánou o diskontní faktor pro interval $(t_0; t_i)$, což by pro opce na akcie je $e^{-r * (t_i - t_0)}$, a $V_i(x)$ označuje hodnotu opce v čase $t_i$, pokud nebyla již vykonána.

  Dle Glassermana (TODO: Zdroj z prezentace, přečíst, aspoň trochu.) můžeme problém stanovení hodnoty opce formulovat takto:

  Nechť $t_m$ označuje čas expirace opce. Potom:
  $$V_m(x) = h_m(x)$$
  $$V_{i-1}(x) = max(h_{i-1}(x), \mathbb(E)\[V_i(X_i) | X_{i-1} = x\])$$

  ...kde $\mathbb(E)$ je označená pro očekávanou hodnotu závorky.

  Nicméně, musíme určit také tzv. pokračovací hodnotu, která udává potenciální zisk z nevyužití opce v tento daný moment. Pokračovací hodnotu budeme označovat $C_i(x)$, a definujeme ji takto:
  $$C_i(x) = \mathbb(E)\[V_{i+1}(X_{i+1} | X_i = x)\]$$

  Nejpoužívanější metodou pro zjištění hodnoty americké opce pomocí Monte Carla je Longstaff-Schwartzova metoda. (TODO: Zdroj prezentace) Mějme $N$ simulovaných cest, počínajících v $t_0$, expirujících v čase $t_m$. V čase expirace je hodnota opce:
  $$ V_m(X_m) = h_m(X_m) $$

  Pokračovací hodnota opce při předchozím datu expirace je:
  $$ C_{m-1}(x) = \mathbb(E)\[V_m(X_m) | X_{m-1} = x\] $$

  Tohle je taková píčovina. Koho to kurva zajímá? Není to o ničem. Stejně nás Trump, Putin, Netanyahu nebo nějakej jinej čurák v moci dostane nukleární pumou, nebo nás dostane nějaká tajná policie. Stejně jsme všichni v prdeli.

  Aproximujeme jako množinu bázových funkcí o velikosti $R$. Formálně psáno jako:
  $$ \hat{C}_{m-1}(x) = \sum\limits^R_{r=1} \beta_r \psi_r (x) $$

  Koeficienty získáme pomocí metody minimalizace rozdílu čtverců. Minimalizujeme postavením derivace rozdílu odhadu od reálné hodnoty rovno nule.

  \begin{equation*}
    \mathbb(E)\left\{ (\mathbb(E)\[V_m(X_m) | X{m-1}\] - \hat{C}_{m-1}(X_{m-1}))^2 \right\}^' = 0
    \mathbb(E)\left\{ (\mathbb(E)\[V_m(X_m) | X{m-1}\] - \hat{C}_{m-1}(X_{m-1})) * \psi_r(X_{m-1}) \right\} = 0
  \end{equation*}

  \begin{align*}
    \mathbb(E) \[V_m(X_m) \psi_r(X_{m-1})\] &= \mathbb(E)\[\hat{C}_{m-1}(X_{m-1} \psi_r (X_{m-1}))\]
    &= \sum\limits_s\mathbb(E)\[\psi_r(X_{m-1})\psi_s(X_{m-1})\]\beta_s
  \end{align*}

